\begingroup
\ArithmeticHead{247}{R}{GAUSS AN SCHUMACHER.}
\SourceLeaf{254}
\pdfbookmark[1]{[7.] Gauss an Schumacher, April 1816}{x1-correspondence-7}
\NumberedItem{7}
\begin{center}\textsc{Gauss} an \textsc{Schumacher}. Göttingen, den\quad\quad April 1816.\end{center}
\HistoricCitation{[Briefwechsel zwischen \textsc{Gauss} und \textsc{Schumacher} I, Altona 1860, S. 125—126.]}
\noindent .\enspace .\enspace .\enspace .\quad Haben Sie denn wirklich vergessen, dass das Arithmetisch Geometrische Mittel, mit welchem Hr. \textsc{Degen} sich beschäftigt, \textls[120]{ganz dasselbe} ist, womit ich mich seit 1791 beschäftigt habe, und jetzt einen ziemlichen Quartband darüber schreiben könnte? Ich habe zwar ausser jenem auch noch andere ar[ithmetisch-]geo[metrische] Mittel betrachtet, die aber ganz elementarisch sind\textsuperscript{[*)]}. Jenes ist das wahre, worüber Sie hier auch eine im Jahre 1800 von mir angefangene kleine Abhandlung gelesen haben (in einem blauen Octavbande, \textit{Varia} betitelt\textsuperscript{[**)]}, worin noch von Ihrer Hand eine Restitutio in Integrum einiger durch einen Dintenfleck unkenntlich gewordenen Stellen zu sehen ist).

\par\indent In jener Abhandlung steht theils ein Beweis, dass wenn ein Ausdruck
\[
\frac{1}{\surd(g+h\cos\varphi)}
\]
\noindent in die Reihe
\[
A+B\cos\varphi+C\cos2\varphi+D\cos3\varphi\ldots
\]
\noindent verwandelt wird, $\dfrac1A$ das Medium Arithm. Geom. zwischen dem grössten und kleinsten Werthe von $\surd(g+h\cos\varphi)$, d. i. zwischen $\surd(g+h)$ und $\surd(g-h)$ ist, theils der Beweis, dass
\[
\begin{aligned}
d\operatorname{Med}(x,y)&=\operatorname{Med}(x,y)\times\\
&\quad\left\{\begin{aligned}
&\frac{dx}{x}\, .\,\left(\frac12+\frac14\, .\,\frac{x-x'}{x'}+\frac18\, .\,\frac{x-x'}{x'}\, .\,\frac{x'-x''}{x''}+\text{etc.}\right)\\[6pt]
{}+&\frac{dy}{y}\, .\,\left(\frac12-\frac14\, .\,\frac{x-x'}{x'}-\frac18\, .\,\frac{x-x'}{x'}\, .\,\frac{x'-x''}{x''}-\text{etc.}\right)
\end{aligned}\right\}
\end{aligned}
\]
\noindent Gerade diese Formel unter andern ist von Ihnen selbst restituirt.

\par\indent Durch die Media und ihre Differentiale kann man dann leicht die übrigen Coefficienten $B,C,D$ etc. bestimmen, so wie überhaupt die Coefficienten von
\par\vfill\noindent\rule{.25\linewidth}{.4pt}\par\smallskip
\begingroup\footnotesize
\noindent [*) Vergl. den Brief [2.] und den zugehörigen \textls[120]{Anhang}, oben S. 234—237.]\par
\noindent [**) Dieser Band befindet sich unter der Bezeichnung »Handbuch 15, Ba« im Gaussarchiv; die betreffende Abhandlung steht S. 3—18, sie ist Werke III, S. 361—374 abgedruckt.]\par\endgroup
\clearpage
\ArithmeticHead{248}{L}{ANALYSIS. BRIEFWECHSEL.}
\SourceLeaf{255}
\noindent $(g+h\cos\varphi)^{\frac12-k}$, wo $k$ irgend eine ganze Zahl ist (welche Reduction bekannt genug ist).

\par\indent In dem zweiten Theile der Abhandlung \textit{Disquisitiones Generales circa Seriem infinitam} $1+\dfrac{\alpha\, .\,\beta}{1\, .\,\gamma}x$ etc. (welche ich vielleicht bald gebe) werde ich einen Theil meiner Untersuchungen über die Ar. Geom. Mittel bekannt zu machen anfangen.

\par\indent Leben Sie wohl, werthester Freund, und vergessen sie nicht
\begin{flushright}Ihren ergebensten\par\medskip C. F. G[\textsc{auss.}]\end{flushright}
\ClosingRule
\clearpage
\endgroup
